\documentclass[12pt, twoside]{article}
\input{../preamble}

\fancyhead[LO]{La Scuola d'Italia / Dr.\ Huson / IB Math \\* 26 March 2026}
\fancyhead[RO]{Markscheme}
\fancyfoot[C]{\thepage}

\begin{document}
\subsubsection*{5.7 Classwork: Analytic Geometry (no calculator) --- Markscheme}

\begin{enumerate}[itemsep=0.2cm]

%--- Problem 1: Line through axes [8 marks] ---
\item A straight line $L$ passes through the points $A(6, 0)$ and $B(0, -4)$.

\begin{enumerate}[itemsep=0.3cm]
  \item Find the gradient of $L$. \hfill [2]

  \begin{minipage}[t]{0.62\textwidth}
  $m = \dfrac{0 - (-4)}{6 - 0} = \dfrac{4}{6}$\\[0.15cm]
  $= \dfrac{2}{3}$
  \end{minipage}%
  \begin{minipage}[t]{0.35\textwidth}\raggedleft
  \textbf{(M1)}\\[0.15cm]
  \textbf{(A1)}
  \end{minipage}

  \medskip
  \item Write down the $y$-intercept of $L$. \hfill [1]

  \begin{minipage}[t]{0.62\textwidth}
  $y$-intercept $= -4$
  \end{minipage}%
  \begin{minipage}[t]{0.35\textwidth}\raggedleft
  \textbf{(A1)}
  \end{minipage}

  \medskip
  \item Find the equation of $L$ in the form $ax + by + d = 0$,
        where $a$, $b$, $d \in \mathbb{Z}$. \hfill [2]

  \begin{minipage}[t]{0.62\textwidth}
  $y = \dfrac{2}{3}x - 4$\\[0.15cm]
  $3y = 2x - 12$\\[0.15cm]
  $2x - 3y - 12 = 0$
  \end{minipage}%
  \begin{minipage}[t]{0.35\textwidth}\raggedleft
  \textbf{(M1)}\\[0.15cm]
  ~\\[0.15cm]
  \textbf{(A1)}
  \end{minipage}

  \medskip
  \item A second line $L_2$ is perpendicular to $L$ and passes through the
        point $A$. Find the equation of $L_2$. \hfill [3]

  \begin{minipage}[t]{0.62\textwidth}
  Perpendicular gradient: $m_2 = -\dfrac{3}{2}$\\[0.2cm]
  $y - 0 = -\dfrac{3}{2}(x - 6)$\\[0.15cm]
  $y = -\dfrac{3}{2}x + 9$
  \end{minipage}%
  \begin{minipage}[t]{0.35\textwidth}\raggedleft
  \textbf{(A1)}\\[0.2cm]
  \textbf{(M1)}\\[0.15cm]
  \textbf{(A1)}
  \end{minipage}
\end{enumerate}

\newpage

%--- Problem 2: Parabola properties [9 marks] ---
\item The graph of $f(x) = x^2 - 6x + 5$ is a parabola.

\begin{enumerate}[itemsep=0.3cm]
  \item Find the $x$-intercepts of the graph of $f$. \hfill [2]

  \begin{minipage}[t]{0.62\textwidth}
  $x^2 - 6x + 5 = 0$\\[0.15cm]
  $(x - 1)(x - 5) = 0$\\[0.15cm]
  $x = 1$ and $x = 5$
  \end{minipage}%
  \begin{minipage}[t]{0.35\textwidth}\raggedleft
  \textbf{(M1)}\\[0.15cm]
  ~\\[0.15cm]
  \textbf{(A1)}
  \end{minipage}

  \medskip
  \item Write $f(x)$ in the form $(x - h)^2 + k$. \hfill [2]

  \begin{minipage}[t]{0.62\textwidth}
  $f(x) = (x^2 - 6x + 9) - 9 + 5$\\[0.15cm]
  $= (x - 3)^2 - 4$
  \end{minipage}%
  \begin{minipage}[t]{0.35\textwidth}\raggedleft
  \textbf{(M1)}\\[0.15cm]
  \textbf{(A1)}
  \end{minipage}

  \medskip
  \item Hence write down the coordinates of the vertex of the parabola. \hfill [1]

  \begin{minipage}[t]{0.62\textwidth}
  Vertex $= (3, -4)$
  \end{minipage}%
  \begin{minipage}[t]{0.35\textwidth}\raggedleft
  \textbf{(A1)}
  \end{minipage}

  \medskip
  \item The graph of $f$ is translated by the vector $\displaystyle\binom{3}{2}$.
        Find the equation of the new graph in the form $y = x^2 + bx + c$. \hfill [4]

  \begin{minipage}[t]{0.62\textwidth}
  New vertex: $(3 + 3,\; -4 + 2) = (6, -2)$\\[0.15cm]
  $y = (x - 6)^2 - 2$\\[0.15cm]
  $= x^2 - 12x + 36 - 2$\\[0.15cm]
  $= x^2 - 12x + 34$
  \end{minipage}%
  \begin{minipage}[t]{0.35\textwidth}\raggedleft
  \textbf{(A1)(A1)}\\[0.15cm]
  \textbf{(M1)}\\[0.15cm]
  ~\\[0.15cm]
  \textbf{(A1)}
  \end{minipage}
\end{enumerate}

\newpage

%--- Problem 3: Line tangent to parabola / discriminant [10 marks] ---
\item The line $y = 2x + k$ is tangent to the parabola $y = x^2 - 4x + 7$.

\begin{enumerate}[itemsep=0.3cm]
  \item Write down an equation whose solutions are the $x$-coordinates
        of the points of intersection of the line and the parabola. \hfill [2]

  \begin{minipage}[t]{0.62\textwidth}
  $2x + k = x^2 - 4x + 7$\\[0.15cm]
  $x^2 - 6x + (7 - k) = 0$
  \end{minipage}%
  \begin{minipage}[t]{0.35\textwidth}\raggedleft
  \textbf{(M1)}\\[0.15cm]
  \textbf{(A1)}
  \end{minipage}

  \medskip
  \item Hence show that $x^2 - 6x + (7 - k) = 0$. \hfill [1]

  \begin{minipage}[t]{0.62\textwidth}
  $x^2 - 4x + 7 - 2x - k = 0$\\[0.15cm]
  $x^2 - 6x + (7 - k) = 0$
  \end{minipage}%
  \begin{minipage}[t]{0.35\textwidth}\raggedleft
  \textbf{(A1)} \textbf{(AG)}
  \end{minipage}

  \medskip
  \item State the condition on the discriminant for the line to be tangent\\
        to the parabola. \hfill [1]

  \begin{minipage}[t]{0.62\textwidth}
  $\Delta = 0$ \quad (the quadratic has exactly one solution)
  \end{minipage}%
  \begin{minipage}[t]{0.35\textwidth}\raggedleft
  \textbf{(A1)}
  \end{minipage}

  \medskip
  \item Find the value of $k$. \hfill [3]

  \begin{minipage}[t]{0.62\textwidth}
  $\Delta = b^2 - 4ac = (-6)^2 - 4(1)(7 - k)$\\[0.15cm]
  $= 36 - 28 + 4k = 8 + 4k$\\[0.15cm]
  $8 + 4k = 0 \Rightarrow k = -2$
  \end{minipage}%
  \begin{minipage}[t]{0.35\textwidth}\raggedleft
  \textbf{(M1)}\\[0.15cm]
  \textbf{(A1)}\\[0.15cm]
  \textbf{(A1)}
  \end{minipage}

  \medskip
  \item Find the coordinates of the point of tangency. \hfill [3]

  \begin{minipage}[t]{0.62\textwidth}
  $x^2 - 6x + 9 = 0$\\[0.15cm]
  $(x - 3)^2 = 0 \Rightarrow x = 3$\\[0.15cm]
  $y = 2(3) + (-2) = 4$\\[0.15cm]
  Point of tangency: $(3, 4)$
  \end{minipage}%
  \begin{minipage}[t]{0.35\textwidth}\raggedleft
  \textbf{(M1)}\\[0.15cm]
  \textbf{(A1)}\\[0.15cm]
  ~\\[0.15cm]
  \textbf{(A1)}
  \end{minipage}
\end{enumerate}

\newpage

%--- Problem 4: Kinematics (pre-calculus) [9 marks] ---
\item A ball is thrown vertically upward. Its height $h$\;metres above the ground
at time $t$ seconds is modelled by
\[h(t) = -5t^2 + 20t + 1, \quad \text{for } t \geq 0.\]

The velocity of the ball at time $t$ is given by the gradient of the
tangent to the curve $h(t)$. For a quadratic $h(t) = at^2 + bt + c$,
the gradient of the tangent at $t$ is $h'(t) = 2at + b$.

\begin{enumerate}[itemsep=0.3cm]
  \item Write down the height of the ball when $t = 0$. \hfill [1]

  \begin{minipage}[t]{0.62\textwidth}
  $h(0) = -5(0) + 20(0) + 1 = 1$\;m
  \end{minipage}%
  \begin{minipage}[t]{0.35\textwidth}\raggedleft
  \textbf{(A1)}
  \end{minipage}

  \medskip
  \item Find $h'(t)$, the velocity of the ball at time $t$. \hfill [2]

  \begin{minipage}[t]{0.62\textwidth}
  $a = -5$, $b = 20$\\[0.15cm]
  $h'(t) = 2(-5)t + 20 = -10t + 20$
  \end{minipage}%
  \begin{minipage}[t]{0.35\textwidth}\raggedleft
  \textbf{(M1)}\\[0.15cm]
  \textbf{(A1)}
  \end{minipage}

  \medskip
  \item Find the time when the ball reaches its maximum height. \hfill [2]

  \begin{minipage}[t]{0.62\textwidth}
  Maximum when $h'(t) = 0$:\\[0.15cm]
  $-10t + 20 = 0 \Rightarrow t = 2$\;s
  \end{minipage}%
  \begin{minipage}[t]{0.35\textwidth}\raggedleft
  \textbf{(M1)}\\[0.15cm]
  \textbf{(A1)}
  \end{minipage}

  \medskip
  \item Find the maximum height. \hfill [2]

  \begin{minipage}[t]{0.62\textwidth}
  $h(2) = -5(4) + 20(2) + 1$\\[0.15cm]
  $= -20 + 40 + 1 = 21$\;m
  \end{minipage}%
  \begin{minipage}[t]{0.35\textwidth}\raggedleft
  \textbf{(M1)}\\[0.15cm]
  \textbf{(A1)}
  \end{minipage}

  \medskip
  \item Find the time when the ball hits the ground. \hfill [2]

  \begin{minipage}[t]{0.62\textwidth}
  $h(t) = 0$: \quad $-5t^2 + 20t + 1 = 0$\\[0.15cm]
  $5t^2 - 20t - 1 = 0$\\[0.15cm]
  $t = \dfrac{20 \pm \sqrt{400 + 20}}{10}
     = \dfrac{20 \pm \sqrt{420}}{10}$\\[0.3cm]
  $t = 2 + \dfrac{\sqrt{105}}{5}$\;s
  \quad ($t > 0$)
  \end{minipage}%
  \begin{minipage}[t]{0.35\textwidth}\raggedleft
  \textbf{(M1)}\\[0.15cm]
  ~\\[0.15cm]
  ~\\[0.3cm]
  \textbf{(A1)}
  \end{minipage}
\end{enumerate}

\end{enumerate}
\end{document}
